\section{Balance armónico y construcción de semillas}
\label{bre:bal:sec}

La localización comienza al separar una parte lineal resoluble y una
no linealidad cuyos coeficientes de Fourier puedan calcularse. El balance
determina frecuencia, amplitud y, cuando es necesaria, una componente
media. Estos datos se convierten en un estado inicial mediante el
resolvente del sistema; después se integra el campo original o una
continuación que termina exactamente en él. La construcción sigue el
método de función descriptiva y localización por continuación
\cite{GelbVanderVelde1968,LeonovKuznetsov2013,KuznetsovEtAl2017}.

\subsection{Realización escalar y convención de transferencia}
\label{bre:bal:lure}

Considérese un sistema autónomo con orden conmensurable
\(0<q\leq1\), terminal inferior fijo \(t_0\) y realización exacta
\begin{equation}
 \Caputo{q}X=PX+b\psi(\sigma),\qquad
 \sigma=r^{\T}X,
 \label{bre:bal:lure-eq}
\end{equation}
donde \(X\in\R^n\), \(P\in\R^{n\times n}\) y
\(b,r\in\R^n\). Para \(q=1\), el operador es la derivada ordinaria.
La palabra exacta exige comprobar
\(PX+b\psi(r^{\T}X)-F(X)\equiv0\) en el dominio del campo.
En particular, una aproximación local de \(F\) no constituye por sí
sola una realización de Lur'e del sistema completo.

Para Chua, escriba \(h(x)=\ell x+\psi(x)\) en
\begin{equation}
 \begin{aligned}
 \Caputo{q}x&=\alpha[y-x-h(x)],\\
 \Caputo{q}y&=x-y+z,\\
 \Caputo{q}z&=-\beta y-\gamma z.
 \end{aligned}
 \label{bre:bal:chua}
\end{equation}
La primera componente se expande como
\(-\alpha(1+\ell)x+\alpha y-\alpha\psi(x)\). Por ello,
\begin{equation}
 \begin{gathered}
 P=\begin{pmatrix}
 -\alpha(1+\ell)&\alpha&0\\
 1&-1&1\\0&-\beta&-\gamma
 \end{pmatrix},\\
 b=(-\alpha,0,0)^{\T},\qquad r=(1,0,0)^{\T}.
 \end{gathered}
 \label{bre:bal:chua-matrices}
\end{equation}
Las dos características consideradas corresponden a
\begin{equation}
 \begin{array}{c|c|c}
 &\ell&\psi(\sigma)\\\hline
 \text{saturación}&m_1&\Delta\operatorname{sat}(\sigma)\\
 \text{arctangente}&a_1&a_2\arctan(\rho\sigma)
 \end{array}
 \label{bre:bal:nonlinearities}
\end{equation}
con \(\Delta=m_0-m_1\), \(\rho>0\) y
\(\operatorname{sat}(s)=\max(-1,\min(1,s))\). El signo negativo de
\(b\) se conserva: cambiarlo sin modificar \(\psi\) cambia el campo.

Tras trasladar \(t_0\) a cero, la transformada de Laplace satisface
\begin{equation}
 (s^qI-P)\widehat X(s)=s^{q-1}X(t_0)+bU(s),
 \label{bre:bal:laplace}
\end{equation}
donde \(U\) es la transformada de \(\psi(\sigma)\). Se utiliza
exclusivamente la convención
\begin{equation}
 \begin{aligned}
 W_q(s)&=r^{\T}(s^qI-P)^{-1}b,\\
 \widehat\sigma(s)
 &=r^{\T}(s^qI-P)^{-1}s^{q-1}X(t_0)\\
 &\quad+W_q(s)U(s).
 \end{aligned}
 \label{bre:bal:transfer}
\end{equation}
La transferencia relaciona entrada y salida después de separar la respuesta
del dato inicial. Las expresiones con resolvente requieren
\(\det(s^qI-P)\ne0\).

\subsection{Resolvente, polinomio y condición de fase de Chua}
\label{bre:bal:resolvent}

Defina \(\lambda=s^q\), \(a_P=\alpha(1+\ell)\),
\(d_0=\beta+\gamma\) y \(d_1=1+\gamma\). Entonces
\begin{equation}
 \lambda I-P=\begin{pmatrix}
 \lambda+a_P&-\alpha&0\\
 -1&\lambda+1&-1\\
 0&\beta&\lambda+\gamma
 \end{pmatrix}.
 \label{bre:bal:resolvent-matrix}
\end{equation}
El menor superior izquierdo es
\begin{equation}
 Q(\lambda)=(\lambda+1)(\lambda+\gamma)+\beta
 =\lambda^2+d_1\lambda+d_0.
 \label{bre:bal:minor}
\end{equation}
La expansión por la primera fila da
\begin{equation}
 \begin{aligned}
 D_\ell(\lambda)
 &:={\det}(\lambda I-P)\\
 &=(\lambda+a_P)Q(\lambda)-\alpha(\lambda+\gamma)\\
 &=\lambda^3+(d_1+a_P)\lambda^2\\
 &\quad +(d_0+a_Pd_1-\alpha)\lambda+a_Pd_0-\alpha\gamma.
 \end{aligned}
 \label{bre:bal:characteristic}
\end{equation}
Al reemplazar la primera columna por \(b\), el determinante de Cramer
es \(-\alpha Q(\lambda)\). Como \(r^{\T}\) selecciona la primera
coordenada,
\begin{equation}
 W_q(s)=-\frac{\alpha Q(s^q)}{D_\ell(s^q)}.
 \label{bre:bal:chua-transfer}
\end{equation}
Este cálculo fija simultáneamente numerador, denominador y signo.

Para \(q=1\), la ecuación de fase se reduce a un polinomio real.
Ponga \(z=\omega^2\), con \(\omega>0\), y escriba
\begin{equation}
 \begin{aligned}
 Q(\ii\omega)&=d_0-z+\ii d_1\omega,\\
 D_\ell(\ii\omega)&=D_R+\ii D_I,\\
 D_R&=a_Pd_0-\alpha\gamma-(d_1+a_P)z,\\
 D_I&=\omega(d_0+a_Pd_1-\alpha-z).
 \end{aligned}
 \label{bre:bal:phase-parts}
\end{equation}
Multiplicar por el conjugado del denominador proporciona
\begin{equation}
 \Im W_1(\ii\omega)
 =-\alpha\frac{d_1\omega D_R-(d_0-z)D_I}{D_R^2+D_I^2}.
 \label{bre:bal:phase-imag}
\end{equation}
Fuera de los polos, dividir el numerador por \(\omega\) y agrupar
términos equivale a resolver
\begin{equation}
 \begin{aligned}
 0={}&z^2+(\alpha+d_1^2-2d_0)z\\
 &+d_0^2-\alpha d_0+\alpha\gamma d_1.
 \end{aligned}
 \label{bre:bal:phase-polynomial}
\end{equation}
Sólo se conservan raíces \(z>0\) que mantengan
\(D_\ell(\ii\sqrt z)\ne0\). Para \(q<1\) se evalúa
\eqref{bre:bal:chua-transfer} en
\begin{equation}
 \zeta=(\ii\omega)^q
 =\omega^q\left(\cos\frac{\pi q}{2}
       +\ii\sin\frac{\pi q}{2}\right),
 \label{bre:bal:spectral-frequency}
\end{equation}
con la rama principal. La frecuencia temporal es \(\omega\);
\(\zeta\) es la variable espectral que entra en el resolvente.

\subsection{Proyección de Fourier y balance centrado}
\label{bre:bal:df}

Para \(\sigma=A\cos\theta\), \(A>0\), el fasor fundamental de la
salida estática es
\begin{equation}
 Y_1(A)=\frac1\pi\int_0^{2\pi}
 \psi(A\cos\theta)e^{-\ii\theta}\,\dd\theta,
 \qquad N(A)=Y_1(A)/A.
 \label{bre:bal:fourier-df}
\end{equation}
La parte seno se anula porque \(\psi(A\cos\theta)\) es par en
\(\theta\). Si además \(\psi\) es impar, su media es cero y, mediante
un desplazamiento de fase y simetría entre semiperiodos,
\begin{equation}
 N(A)=\frac{2}{\pi A}\int_0^\pi
 \psi(A\sin\theta)\sin\theta\,\dd\theta\in\R.
 \label{bre:bal:df-sine}
\end{equation}
La integral pertenece a la característica estática: no contiene \(q\).
El orden fraccionario entra a través del bloque dinámico \(W_q\).
Si la no linealidad posee histéresis o estado interno, debe calcularse el
fasor de su respuesta periódica; en ese caso \(N\) puede ser complejo
y depender de la frecuencia y de la historia interna.

Al retener sólo el fundamental, la salida del bloque lineal debe
reproducir la entrada \(\sigma\). Así,
\(A=W_q(\ii\omega)Y_1\), y la ecuación de cierre es
\begin{equation}
 1-W_q(\ii\omega)N(A)=0.
 \label{bre:bal:physical-closure}
\end{equation}
Para \(N\) real, una solución finita no nula exige
\begin{equation}
 \Im W_q(\ii\omega)=0,\qquad
 N(A)=k:=\frac1{\Re W_q(\ii\omega)}.
 \label{bre:bal:phase-gain}
\end{equation}
Después de la condición de fase debe comprobarse que \(k\) pertenece al
rango real de \(N\). Una intersección de fase cuya ganancia no puede
alcanzar la no linealidad no produce amplitud admisible.

\subsubsection{Saturación}
Para \(\psi(s)=\Delta\operatorname{sat}(s)\) y \(0<A\leq1\),
\(\psi(A\sin\theta)=\Delta A\sin\theta\); como
\(\int_0^\pi\sin^2\theta\,\dd\theta=\pi/2\), resulta
\(N(A)=\Delta\). Si \(A>1\), sea
\(\theta_0=\arcsin(1/A)\). Los tramos no saturados son
\([0,\theta_0]\) y \([\pi-\theta_0,\pi]\); en el tramo central la
saturación vale uno. Por ello,
\begin{equation}
 \begin{aligned}
 N(A)&=\frac{2\Delta}{\pi A}
 \left[2A\int_0^{\theta_0}\sin^2\theta\,\dd\theta
       +2\cos\theta_0\right]\\
 &=\frac{2\Delta}{\pi A}
 \left[A\theta_0-\frac A2\sin(2\theta_0)+2\cos\theta_0\right]\\
 &=\frac{2\Delta}{\pi}
 \left[\arcsin\frac1A+\frac{\sqrt{A^2-1}}{A^2}\right].
 \end{aligned}
 \label{bre:bal:saturation-df}
\end{equation}
El último paso usa \(A\sin\theta_0=1\). Para \(\Delta>0\),
\begin{equation}
 N'(A)=-\frac{4\Delta\sqrt{A^2-1}}{\pi A^3}<0
 \quad(A>1).
 \label{bre:bal:saturation-derivative}
\end{equation}
Además, \(N(1)=\Delta\) y \(N(A)\to0\) al crecer \(A\).
Existe una única amplitud saturada para \(0<k<\Delta\). Si
\(k=\Delta\), toda amplitud \(0<A\leq1\) satisface el balance lineal
y éste no selecciona una amplitud aislada.

\subsubsection{Arctangente}
Sea \(\psi(s)=a_2\arctan(\rho s)\), \(\xi=\rho A\), y defina
\(I(\xi)=\int_0^\pi\arctan(\xi\sin\theta)\sin\theta\,\dd\theta\).
La diferenciación bajo la integral es válida por continuidad y acotación
del integrando y de su derivada en intervalos compactos de \(\xi\).
Para \(\xi>0\),
\begin{equation}
 \begin{aligned}
 I'(\xi)&=\int_0^\pi
 \frac{\sin^2\theta}{1+\xi^2\sin^2\theta}\,\dd\theta\\
 &=\frac1{\xi^2}\left[\pi-
 \int_0^\pi\frac{\dd\theta}{1+\xi^2\sin^2\theta}\right].
 \end{aligned}
 \label{bre:bal:arctan-integral}
\end{equation}
La integral restante se reduce por simetría a dos veces el intervalo
\([0,\pi/2]\). Con \(u=\tan\theta\), se transforma en
\(2\int_0^\infty[1+(1+\xi^2)u^2]^{-1}\,\dd u
=\pi/\sqrt{1+\xi^2}\). La función
\(\pi(\sqrt{1+\xi^2}-1)/\xi\) tiene la derivada obtenida y límite
cero en \(\xi=0\), igual a \(I(0)\). Se concluye
\begin{equation}
 \begin{aligned}
 N(A)&=\frac{2a_2}{\pi A}I(\rho A)\\
 &=\frac{2a_2}{\rho A^2}
       \left(\sqrt{1+\rho^2A^2}-1\right)\\
 &=\frac{2a_2\rho}{1+\sqrt{1+\rho^2A^2}}.
 \end{aligned}
 \label{bre:bal:arctan-df}
\end{equation}
La última forma evita la resta de números próximos cuando \(A\) es
pequeño. Si \(a_2\ne0\), la ecuación \(N(A)=k\) se despeja como
\begin{equation}
 A^2=\frac{4a_2(a_2\rho-k)}{k^2\rho}.
 \label{bre:bal:arctan-amplitude}
\end{equation}
Antes de elevar al cuadrado se tiene
\(\sqrt{1+\rho^2A^2}=2a_2\rho/k-1\). Por tanto, para \(A>0\),
la condición completa es \(0<k/(a_2\rho)<1\); impone el signo
correcto y evita aceptar raíces introducidas por el cuadrado.

\subsection{Balance de media y fundamental sesgado}
\label{bre:bal:bias}

Para localizar una oscilación alrededor de una media distinta de cero,
tome \(\sigma=c+A\cos\theta\). Defina los coeficientes físicos
\begin{equation}
 \begin{aligned}
 \Psi_0(c,A)&=\frac1{2\pi}\int_0^{2\pi}
       \psi(c+A\cos\theta)\,\dd\theta,\\
 \Psi_1(c,A)&=\frac1\pi\int_0^{2\pi}
       \psi(c+A\cos\theta)\cos\theta\,\dd\theta,\\
 N_1(c,A)&=\Psi_1(c,A)/A.
 \end{aligned}
 \label{bre:bal:bias-fourier}
\end{equation}
La media del bloque dinámico satisface \(P\bar X+b\Psi_0=0\).
Si \(P\) es invertible,
\begin{equation}
 \bar X=-P^{-1}b\Psi_0,\qquad
 c=r^{\T}\bar X=W_q(0)\Psi_0.
 \label{bre:bal:dc}
\end{equation}
Para Chua,
\begin{equation}
 W_q(0)=-\frac{\beta+\gamma}
 {(1+\ell)(\beta+\gamma)-\gamma}.
 \label{bre:bal:chua-dc}
\end{equation}
Si \(P\) es singular, se resuelven directamente las ecuaciones de
media, sin usar \(P^{-1}\) ni una ganancia continua infinita.

La condición armónica sigue siendo
\begin{equation}
 1-W_q(\ii\omega)N_1(c,A)=0.
 \label{bre:bal:bias-closure}
\end{equation}
La ecuación de media y las partes real e imaginaria de
\eqref{bre:bal:bias-closure} son tres ecuaciones reales para
\((c,A,\omega)\). Mantener la media evita tratar el desplazamiento
del ciclo como una coordenada libre sin ecuación.

Para la saturación, las integrales se evalúan por tramos incluso cuando
la entrada queda totalmente saturada. Con
\([u]_{-1}^{1}=\max(-1,\min(1,u))\), sean
\begin{equation}
 \theta_+=\arccos\left[\frac{1-c}{A}\right]_{-1}^{1},\qquad
 \theta_-=\arccos\left[\frac{-1-c}{A}\right]_{-1}^{1}.
 \label{bre:bal:bias-crossings}
\end{equation}
Se cumple \(0\leq\theta_+\leq\theta_-\leq\pi\). En
\([0,\pi]\), la salida normalizada vale sucesivamente
\(1\), \(c+A\cos\theta\) y \(-1\), con cambios en esos
ángulos. Integrar cada tramo da
\begin{equation}
 \begin{aligned}
 \Psi_0=\frac\Delta\pi\bigl[&\theta_++\theta_--\pi
       +c(\theta_--\theta_+)\\
       &+A(\sin\theta_--\sin\theta_+)\bigr].
 \end{aligned}
 \label{bre:bal:bias-sat-dc}
\end{equation}
Para el fundamental se usan las primitivas
\(\int\cos\theta\,\dd\theta=\sin\theta\) y
\(\int\cos^2\theta\,\dd\theta=\theta/2+\sin(2\theta)/4\).
La suma de los tres tramos es
\begin{equation}
 \begin{aligned}
 \Psi_1=\frac{2\Delta}\pi\bigl[&\sin\theta_++\sin\theta_-
       +c(\sin\theta_--\sin\theta_+)\\
       &+\tfrac A2(\theta_--\theta_+)\\
       &+\tfrac A4(\sin2\theta_--\sin2\theta_+)\bigr].
 \end{aligned}
 \label{bre:bal:bias-sat-first}
\end{equation}
Los ángulos recortados hacen desaparecer automáticamente un tramo de
longitud cero. En \(c=0\), la media se anula y
\(\Psi_1/A\) recupera \eqref{bre:bal:saturation-df}.

\subsection{Modo normalizado, semilla y equivalencias}
\label{bre:bal:modal}

Fijada una raíz de balance, sea \(k=N(A)\) o \(k=N_1(c,A)\), y
\(P_k=P+bkr^{\T}\). Si \(\zeta I-P\) es invertible, el lema del
determinante matricial proporciona
\begin{equation}
 \begin{aligned}
 \det(\zeta I-P_k)
 &=\det(\zeta I-P)\\
 &\quad\times\left[1-kr^{\T}(\zeta I-P)^{-1}b\right].
 \end{aligned}
 \label{bre:bal:det-lemma}
\end{equation}
En \(\zeta=(\ii\omega)^q\), el corchete es cero por el cierre.
Un modo normalizado se construye directamente, sin una elección
arbitraria de escala, mediante
\begin{equation}
 v=k(\zeta I-P)^{-1}b.
 \label{bre:bal:mode}
\end{equation}
En efecto, \(r^{\T}v=kW_q(\ii\omega)=1\) y
\((\zeta I-P_k)v=bk-bkr^{\T}v=0\).
Para Chua, las dos últimas filas de esta ecuación dan
\begin{equation}
 v=\left(1,\frac{\zeta+\gamma}{Q(\zeta)},
               -\frac\beta{Q(\zeta)}\right)^{\T}.
 \label{bre:bal:chua-mode}
\end{equation}
La igualdad presupone \(Q(\zeta)\ne0\), condición que se comprueba
junto con la existencia del resolvente.

En la rama sesgada, \(\Psi_1=kA\) y
\(X_1=(\zeta I-P)^{-1}b\Psi_1=Av\). Se obtiene la familia de semillas
\begin{equation}
 X_{\mathrm{sem}}(\varphi)=\bar X+Re(Av e^{\ii\varphi}),
 \qquad 0\leq\varphi<2\pi,
 \label{bre:bal:seed}
\end{equation}
con \(\bar X=0\) en la rama centrada. La salida es exactamente
\(r^{\T}X_{\mathrm{sem}}=c+A\cos\varphi\). El par
\(\zeta,\overline\zeta\) está en la frontera angular
\(|\arg\zeta|=q\pi/2\) del auxiliar. Para usarlo como predictor de
una oscilación dominante se comprueba que los demás modos satisfagan el
criterio angular estricto de estabilidad; el balance por sí solo no
establece esa propiedad \cite{Matignon1996}.

Una similitud real \(X=SY\), con \(S\) invertible, transforma
\(P,b,r^{\T}\) en \(S^{-1}PS,S^{-1}b,r^{\T}S\).
Como
\((s^qI-S^{-1}PS)^{-1}=S^{-1}(s^qI-P)^{-1}S\),
la transferencia no cambia y la semilla se transforma por
\(Y_{\mathrm{sem}}=S^{-1}X_{\mathrm{sem}}\).
Si se añade un desplazamiento \(X=SY+d\), deben conservarse también
el término constante \(S^{-1}Pd\) y el argumento
\(r^{\T}SY+r^{\T}d\); si \(d\) es equilibrio, se cancelan usando
\(Pd+b\psi(r^{\T}d)=0\) y la no linealidad incremental
\(\psi(r^{\T}SY+r^{\T}d)-\psi(r^{\T}d)\).

Otra equivalencia consiste en redistribuir una ganancia \(a>0\):
\(\widetilde b=ab\), \(\widetilde\psi=\psi/a\).
Entonces \(\widetilde W_q=aW_q\),
\(\widetilde N=N/a\), y
\(\widetilde W_q\widetilde N=W_qN\). El campo, el cierre y el modo
\(\widetilde k(\zeta I-P)^{-1}\widetilde b=v\) permanecen iguales.
Esta comprobación permite comparar realizaciones con diferentes escalas
sin confundir una ganancia de representación con un cambio del modelo.

\subsection{Del armónico estacionario al PVI causal}
\label{bre:bal:weyl-caputo}

El balance usa la acción del operador estacionario de Liouville--Weyl
sobre los armónicos:
\(D_{\mathrm W}^{q}e^{\ii\omega t}=(\ii\omega)^qe^{\ii\omega t}\).
Para \(0<q<1\), esta identidad se interpreta por Fourier o mediante
regularización de Abel de la integral sobre toda la historia. En efecto,
para \(\varepsilon>0\),
\begin{equation}
 \int_0^\infty u^{-q}e^{-(\varepsilon+\ii\omega)u}\,\dd u
 =\Gamma(1-q)(\varepsilon+\ii\omega)^{q-1}.
 \label{bre:bal:abel}
\end{equation}
Multiplicar por \(\ii\omega/\Gamma(1-q)\) y tomar
\(\varepsilon\downarrow0\) produce la rama de
\eqref{bre:bal:spectral-frequency}.

En el cálculo siguiente se fija \(t_0=0\). La derivada de Caputo
retiene sólo el intervalo \([0,t]\).
Para el mismo armónico suave,
\begin{equation}
 \Caputo{q}e^{\ii\omega t}
 =\frac{\ii\omega e^{\ii\omega t}}{\Gamma(1-q)}
   \int_0^t u^{-q}e^{-\ii\omega u}\,\dd u.
 \label{bre:bal:caputo-harmonic}
\end{equation}
Restar la integral sobre \([0,\infty)\) revela el defecto inicial:
\begin{equation}
 \begin{aligned}
 \Caputo{q}e^{\ii\omega t}-\zeta e^{\ii\omega t}
 ={}&-\frac{\ii\omega e^{\ii\omega t}}{\Gamma(1-q)}\\
 &\times\int_t^\infty u^{-q}e^{-\ii\omega u}\,\dd u.
 \end{aligned}
 \label{bre:bal:initial-defect}
\end{equation}
La integral impropia converge para \(\omega>0\) por el criterio de
Dirichlet. Cerca de cero,
\(\Caputo{q}e^{\ii\omega t}
=\ii\omega t^{1-q}/\Gamma(2-q)+O(t^{2-q})\), mientras
\(\zeta e^{\ii\omega t}\to\zeta\ne0\).
Por consiguiente, incluso el armónico del auxiliar lineal exige una
respuesta de arranque en el PVI de Caputo. Para \(q=1\), la identidad
ordinaria no presenta este defecto.

En el sistema no lineal se añade un segundo defecto: los armónicos
descartados de \(\psi\). El balance de primer armónico no los hace
cero. La semilla de \eqref{bre:bal:seed} localiza un estado, pero la
trayectoria causal debe obtenerse resolviendo el PVI. Los resultados
de no periodicidad para sistemas autónomos de Caputo con terminal
inferior finito impiden interpretar sin más el cierre estacionario como
una órbita periódica exacta del PVI
\cite{TavazoeiHaeri2009,AreaLosadaNieto2014}.

La distinción permite precisar el traslado de los métodos publicados.
Danca utiliza el PVI de Caputo y ABM, y describe la localización mediante
linealización armónica y continuación, junto con exploraciones numéricas
de las vecindades de equilibrios \cite{Danca2017}. Wu et al. aplican a
Chua con arctangente el algoritmo de localización de Leonov y construyen
un auxiliar mediante transferencias racionales en una variable compleja
\(p\), con un bloque de denominador \(p^2+\omega_0^2\)
\cite{Wu2023ArctanChua}. La reducción algebraica de una matriz por
similitud sigue siendo válida en Caputo; interpretar sus polos como una
frecuencia temporal requiere además la relación \(p=(\ii\omega)^q\).
En particular, conservar autovalores \(\pm\ii\omega_0\) con \(q<1\)
los sitúa estrictamente en el sector estable de Matignon, no en su
frontera. Es posible conservar una semilla obtenida en orden entero y
probarla en Caputo, pero eso constituye un PVI distinto del predictor
que recalcula la fase con \(s^q\).

Wu et al. reconocen expresamente tanto la no periodicidad exacta como
la pérdida de parte de la memoria histórica al emplear su aproximación
por descomposición de Adomian en subintervalos. También distinguen los
contornos numéricamente periódicos de la periodicidad matemática
\cite{Wu2023ArctanChua}. El problema pendiente al transferir una
prueba entera es demostrar persistencia y atracción con la memoria
causal declarada; no se resuelve sólo reproduciendo un contorno o
reutilizando una fórmula de estado inicial. La independencia de \(q\)
en \eqref{bre:bal:df-sine} sí es legítima para una característica
estática: elevar \(N\) a \(q\) no se deduce ni de Fourier ni de
la transformada de Caputo.

\subsection{Continuación que conserva la media y el campo final}
\label{bre:bal:continuation}

La descomposición
\(\psi(\sigma)=\Psi_0+k(\sigma-c)
+[\psi(\sigma)-\Psi_0-k(\sigma-c)]\) define el auxiliar afín
\(P_k=P+bkr^{\T}\), \(d_k=b(\Psi_0-kc)\), y la familia
\begin{equation}
 \begin{aligned}
 F_\varepsilon(X)&=P_kX+d_k+\varepsilon b\delta\psi(r^{\T}X),\\
 \delta\psi(\sigma)&=\psi(\sigma)-\Psi_0-k(\sigma-c),
 \qquad 0\leq\varepsilon\leq1.
 \end{aligned}
 \label{bre:bal:homotopy}
\end{equation}
En \(\varepsilon=0\),
\(F_0(\bar X)=P\bar X+bk c+b(\Psi_0-kc)=0\), y
\(F_0(\bar X+\xi)=P_k\xi\). En \(\varepsilon=1\), los términos
\(bk r^{\T}X\), \(b\Psi_0\) y \(bkc\) se cancelan con los del
corchete y queda exactamente \(F_1(X)=PX+b\psi(r^{\T}X)\).
Una homotopía que omita \(d_k\) pierde esta correspondencia cuando
\(c\ne0\).

Se inicia en \(X_{\mathrm{sem}}\), se hace crecer
\(\varepsilon(t)\) por etapas y se evalúa una única integral causal:
\begin{equation}
 \begin{aligned}
 X(t)=X_{\mathrm{sem}}+\frac1{\Gamma(q)}
 \int_{t_0}^t(t-s)^{q-1}F_{\varepsilon(s)}(X(s))\,\dd s.
 \end{aligned}
 \label{bre:bal:causal-continuation}
\end{equation}
Cada valor pasado del campo conserva el parámetro que tuvo en ese
instante. Al cambiar de etapa se prolonga esa historia. La continuidad
numérica de una rama sólo permite seguir candidatos: no demuestra que
exista un atractor en todos los parámetros intermedios. Para un estudio
de cuencas con reinicios se realiza después un PVI autónomo del campo
final con el mismo terminal inferior y convenio de memoria usados en
los demás lanzamientos.

\subsubsection{Recuperación del método entero}
Las matrices \(P,b,r\), las proyecciones \(\Psi_0,\Psi_1\) y
la identidad algebraica de la homotopía no dependen del orden. Para
una frecuencia positiva fija y fuera de los polos,
\begin{equation}
 \begin{aligned}
 (\ii\omega)^q&\longrightarrow\ii\omega,\\
 W_q(\ii\omega)&\longrightarrow r^{\T}(\ii\omega I-P)^{-1}b
 \qquad(q\to1^-).
 \end{aligned}
 \label{bre:bal:integer-limit}
\end{equation}
La continuidad del inverso matricial se obtiene, por ejemplo, de su
expresión mediante la adjunta y el determinante no nulo. Por tanto,
el cierre converge a \(1-W_1N=0\); para Chua, su ecuación de fase
es precisamente \eqref{bre:bal:phase-polynomial}.

La continuidad de las \emph{raíces} requiere una hipótesis adicional.
Para la rama sesgada, forme el residual real
\begin{equation}
 \mathcal F(c,A,\omega;q)=\begin{pmatrix}
 c-W_q(0)\Psi_0\\
 \Re(1-W_q(\ii\omega)N_1)\\
 \Im(1-W_q(\ii\omega)N_1)
 \end{pmatrix}.
 \label{bre:bal:branch-residual}
\end{equation}
Si \(\mathcal F\) es \(C^1\) cerca de una raíz con \(q=1\),
\(A>0\), \(\omega>0\), y su Jacobiano respecto de
\((c,A,\omega)\) es invertible, el teorema de la función implícita
produce una rama local \((c(q),A(q),\omega(q))\) que converge a
esa raíz. Para el balance centrado se aplica el mismo argumento a
las dos ecuaciones reales en \((A,\omega)\).

En esa rama, \(k(q)=N_1(c(q),A(q))\), \(P_{k(q)}\), \(v(q)\),
\(\bar X(q)\) y la semilla de \eqref{bre:bal:seed} convergen a sus
expresiones enteras porque se construyen mediante funciones continuas
y resolventes no singulares. En \(q=1\),
\(\bar X+\Re(Av e^{\ii\omega t})\) resuelve exactamente el auxiliar
afín: al derivarla, \(\ii\omega v=P_kv\) y \(F_0(\bar X)=0\)
dan la igualdad componente a componente. El campo final sigue
requiriendo la continuación y su validación dinámica. Si falla la
invertibilidad del Jacobiano del balance, pueden fusionarse o terminar
ramas; el límite de las fórmulas no demuestra su continuación global.

\subsection{Ejemplo calculado: Chua con saturación y \texorpdfstring{\(q=0.9998\)}{q=0.9998}}
\label{bre:bal:example}

Tome
\begin{equation}
 \begin{gathered}
 (\alpha,\beta,\gamma)=(8.4562,12.0732,0.0052),\\
 (m_0,m_1)=(-0.1768,-1.1468),\qquad q=0.9998.
 \end{gathered}
 \label{bre:bal:example-parameters}
\end{equation}
La separación da \(\ell=-1.1468\), \(\Delta=0.9700\),
\(d_0=12.0784\), \(d_1=1.0052\). Como control entero, el polinomio
de fase de \eqref{bre:bal:phase-polynomial} es
\(z^2-14.69017296z+43.794581375648\). Para el orden fraccionario
indicado, se resuelve directamente \(\Im W_q(\ii\omega)=0\).
Una de las ramas proporciona
\begin{equation}
 \begin{aligned}
 \omega&\simeq2.04028605108,\\
 W_q(\ii\omega)&\simeq4.76138797078,\\
 k&\simeq0.210022792962,\\
 \zeta&\simeq0.000640883348+2.03999498963\ii.
 \end{aligned}
 \label{bre:bal:example-frequency}
\end{equation}
Como \(0<k<\Delta\), la ecuación centrada tiene una amplitud única
saturada, \(A\simeq5.85176778549\). Para localizar una rama sesgada
se conserva la misma condición de fase y se resuelve
\begin{equation}
 c-6.79206998612\Psi_0(c,A)=0,\qquad
 \Psi_1(c,A)-kA=0.
 \label{bre:bal:example-system}
\end{equation}
Las fórmulas por tramos dan, con cifras redondeadas,
\begin{equation}
 \begin{aligned}
 c&\simeq2.776397,& A&\simeq4.577883,\\
 \theta_+&\simeq1.969299,&\theta_-&\simeq2.540861,\\
 \Psi_0&\simeq0.408770354,&\Psi_1&\simeq0.961459759.
 \end{aligned}
 \label{bre:bal:example-bias}
\end{equation}
Se verifica \(W_q(0)\Psi_0\simeq c\) y
\(\Psi_1/A\simeq k\). Las dos últimas ecuaciones de la media dan
\(\bar y=\gamma c/(\beta+\gamma)\) y
\(\bar z=-\beta c/(\beta+\gamma)\), por lo que
\begin{equation}
 \bar X\simeq(2.776397,\ 0.001195296,\ -2.775202)^{\T}.
 \label{bre:bal:example-mean}
\end{equation}
La sustitución de \(\zeta\) en \eqref{bre:bal:chua-mode} produce
\begin{equation}
 v\simeq\begin{pmatrix}
 1\\0.063298582+0.241242519\ii\\
 -1.428794382+0.370525917\ii
 \end{pmatrix}.
 \label{bre:bal:example-mode}
\end{equation}
Para fase \(\varphi=0\), la suma \(\bar X+A\Re v\) proporciona
\begin{equation}
 X_{\mathrm{sem}}\simeq
 (7.354280,\ 0.290969,\ -9.316055)^{\T}.
 \label{bre:bal:example-seed}
\end{equation}
El tercer valor propio del auxiliar es aproximadamente
\(-1.54110635\): se obtiene restando \(2\Re\zeta\) de
\(\operatorname{tr}P_k\), y está en el sector estable. Así quedan
determinados el modo, la media, la escala y el estado inicial. La
continuación usa \eqref{bre:bal:homotopy}; la naturaleza del destino se
establece después mediante integración y análisis de cuencas.

\subsection{Familia de Machado y residuo del modelo físico}
\label{bre:bal:machado}

Machado propone elevar una función descriptiva a un orden real
\cite{Machado2015FDF}. Para distinguir esta operación del orden
dinámico \(q\), se denota su exponente por \(\mu>0\). Con una escala
positiva \(N_\star\) de las mismas unidades que \(N\), la familia es
\begin{equation}
 N_{\mu,\star}=N_\star
 \exp\left[\mu\operatorname{Log}(N/N_\star)\right].
 \label{bre:bal:machado-definition}
\end{equation}
Sobre \(N>0\) es real y unívoca. Para \(N\) complejo se fija una
hoja del logaritmo y se continúa su fase; para \(N<0\), la potencia
principal es, en general, compleja. En \(N=0\), la extensión
\(0^\mu=0\) no permite utilizar el logaritmo. La normalización evita
elevar directamente una magnitud dimensional a un exponente distinto
de uno. No hay una identidad que obligue a imponer \(\mu=q\).

En la rama positiva, el cierre auxiliar
\(1-W_q(\ii\omega)N_{\mu,\star}=0\) conserva la condición de
fase, pero sustituye la ecuación de amplitud por
\begin{equation}
 N(A)=N_\star\left(\frac{k}{N_\star}\right)^{1/\mu},
 \qquad k>0.
 \label{bre:bal:machado-amplitude}
\end{equation}
Se comprueba nuevamente el rango de \(N\). Al derivar
\(\mu\log[N(A)/N_\star]=\log(k/N_\star)\), con \(k\) fijo,
se obtiene
\begin{equation}
 \frac{\dd A}{\dd\mu}
 =-\frac{N(A)\log[N(A)/N_\star]}{\mu N'(A)}.
 \label{bre:bal:machado-sensitivity}
\end{equation}
Esta fórmula identifica pérdida de condicionamiento cuando
\(N'(A)\) se aproxima a cero. El modo se calcula con la ganancia
auxiliar \(k\), y la nueva amplitud modifica la semilla.

La potencia no es el coeficiente de Fourier de la característica
original. Por ello se calculan separadamente
\begin{equation}
 \begin{aligned}
 \rho_{\mathrm{aux}}&=|1-W_q(\ii\omega)N_{\mu,\star}(A)|,\\
 \rho_{\mathrm{fis}}&=|1-W_q(\ii\omega)N(A)|.
 \end{aligned}
 \label{bre:bal:machado-residuals}
\end{equation}
En el ejemplo anterior, con \(N_\star=1\) y \(\mu=2\),
\(A\simeq2.62843351478\), \(N(A)\simeq0.458282437981\) y
\(N(A)^2\simeq k\). El cierre auxiliar tiene residuo del orden de
\(10^{-13}\), pero el físico vale aproximadamente \(1.18206048743\).
La nueva semilla diversifica la búsqueda; no resuelve el balance físico
de primer armónico.

La escala también interviene en la invariancia de representación. Bajo
\(b\mapsto ab\), \(N\mapsto N/a\), debe hacerse
\(N_\star\mapsto N_\star/a\); entonces
\(N_{\mu,\star}\mapsto N_{\mu,\star}/a\), y el cierre auxiliar no
cambia. Si se mantiene una potencia desnuda \(N^\mu\), el producto
se transforma en \(a^{1-\mu}W_qN^\mu\), que depende de la escala
elegida cuando \(\mu\ne1\).

En una rama sesgada se mantiene la media física \(\Psi_0\) y sólo
se aplica la potencia a \(N_1\). Se resuelven
\(c=W_q(0)\Psi_0\) y
\(1-W_qN_{1,\mu,\star}=0\), registrando además
\(|1-W_qN_1|\). La ganancia auxiliar correspondiente se inserta en
la homotopía afín para conservar \(F_1=F\). Esta construcción escalar
no autoriza potencias elemento a elemento en un problema multicanal.

\subsection{Balance multicanal y control del residuo completo}
\label{bre:bal:mimo}

Para no linealidades acopladas se emplea la identidad exacta
\begin{equation}
 F(X)=AX+B\Phi(CX),\qquad
 W_q(s)=C(s^qI-A)^{-1}B,
 \label{bre:bal:mimo-model}
\end{equation}
con \(A\in\R^{n\times n}\), \(B\in\R^{n\times m}\),
\(C\in\R^{p\times n}\) y \(\Phi:\R^p\to\R^m\).
Si \(K=D\Phi(C\bar X)\), la linealización tiene matriz
\(A+BKC\). Fuera de los polos, la identidad
\(\det(I-UV)=\det(I-VU)\) permite escribir
\begin{equation}
 \begin{aligned}
 \det(\zeta I-A-BKC)
 &=\det(\zeta I-A)\\
 &\quad\times\det[I_p-W_q(\ii\omega)K].
 \end{aligned}
 \label{bre:bal:mimo-linear}
\end{equation}
Aquí se fija $\zeta=(\ii\omega)^q$. Este predictor usa una derivada local.
Para amplitudes finitas se
calculan los coeficientes de Fourier de \(\Phi\) sobre la señal
propuesta, sin sustituir sus productos por ganancias escalares separadas.

Sea
\begin{equation}
 X^{(H)}(\theta)=\sum_{k=-H}^{H}X_ke^{\ii k\theta},\qquad
 X_{-k}=\overline{X_k},\quad X_0\in\R^n.
 \label{bre:bal:fourier-state}
\end{equation}
Aquí \(X_1\) es un coeficiente complejo de Fourier: una componente
\(A\cos\theta\) tiene coeficientes \(A/2\) en \(k=\pm1\).
En \eqref{bre:bal:seed}, en cambio, \(Av\) era el fasor completo.
La reconstrucción que conserva esta normalización es
\begin{equation}
 X_{\mathrm{sem}}(\varphi)=X_0+
       2\Re\sum_{k=1}^{H}X_ke^{\ii k\varphi}.
 \label{bre:bal:mimo-seed}
\end{equation}
Si \(u=\sum U_ke^{\ii k\theta}\) y
\(v=\sum V_ke^{\ii k\theta}\), al multiplicar las series y reunir
los términos con el mismo exponente se obtiene
\begin{equation}
 \widehat{uv}_k=\sum_jU_jV_{k-j},\qquad
 \widehat{uvw}_k=\sum_j\sum_lU_jV_lW_{k-j-l}.
 \label{bre:bal:convolutions}
\end{equation}
Las sumas se extienden sobre enteros, tomando coeficientes cero fuera
del intervalo retenido. Para polinomios cuadráticos o cúbicos esto
calcula exactamente la no linealidad de la señal truncada, incluidos
los armónicos generados hasta \(2H\) o \(3H\).

Por ejemplo, para el Lorenz generalizado,
\begin{equation}
 \begin{gathered}
 A=\begin{pmatrix}-\sigma&\sigma&0\\r&-1&0\\0&0&-1\end{pmatrix},
 \quad B=\operatorname{diag}(-a,-1,1),\quad C=I,\\
 \Phi(X)=(yz,xz,xy)^{\T}.
 \end{gathered}
 \label{bre:bal:lorenz-mimo}
\end{equation}
La primera componente del coeficiente no lineal es
\(\widehat\Phi_{1,k}=\sum_jX_{2,j}X_{3,k-j}\); las otras dos se
obtienen de \(xz\) y \(xy\). Para una señal de un solo armónico,
\(x=c_x+a_x\cos\theta\), \(y=c_y+a_y\cos\theta\), el producto
contiene media \(c_xc_y+a_xa_y/2\), fundamental
\((c_xa_y+c_ya_x)\cos\theta\) y segundo armónico
\((a_xa_y/2)\cos2\theta\). Así se ve por qué deben conservarse
tanto la media como el acoplamiento entre canales.

Defina
\(\widehat\Phi_k=(2\pi)^{-1}\int_0^{2\pi}
\Phi(CX^{(H)}(\theta))e^{-\ii k\theta}\,\dd\theta\).
El balance estacionario es
\begin{equation}
 \begin{aligned}
 AX_0+B\widehat\Phi_0&=0,\\
 [ (\ii k\omega)^qI-A]X_k-B\widehat\Phi_k&=0,
 \quad 1\leq k\leq H.
 \end{aligned}
 \label{bre:bal:mimo-balance}
\end{equation}
Si \(X_1\) tiene alguna componente no nula, se añade una condición de fase,
\(\Im(e_j^{\T}X_1)=0\), con \(\Re(e_j^{\T}X_1)>0\).
Si el primer coeficiente se anula, se fija la fase mediante otro
armónico no nulo.
Son \(n+2nH+1\) incógnitas reales, incluida \(\omega\), y el mismo
número de ecuaciones. La desigualdad excluye la solución estacionaria
de amplitud cero en la coordenada elegida.

Resolver las ecuaciones retenidas no anula el residuo fuera de ellas.
Sea
\(\mathcal R(\theta)=D_{\mathrm W}^qX^{(H)}-AX^{(H)}
-B\Phi(CX^{(H)})\). Sus coeficientes son
\begin{equation}
 \begin{aligned}
 \mathcal R_0&=-AX_0-B\widehat\Phi_0,\\
 \mathcal R_k&=[(\ii k\omega)^qI-A]X_k-B\widehat\Phi_k
       \quad(1\leq k\leq H),\\
 \mathcal R_k&=-B\widehat\Phi_k\quad(k>H).
 \end{aligned}
 \label{bre:bal:mimo-residual}
\end{equation}
Para señales reales, los índices negativos son los conjugados. Si el
residuo pertenece a \(L^2\), Parseval, con la norma euclídea, da
\begin{equation}
 \frac1{2\pi}\int_0^{2\pi}\norm{\mathcal R(\theta)}^2\,\dd\theta
 =\norm{\mathcal R_0}^2+2\sum_{k=1}^{\infty}\norm{\mathcal R_k}^2.
 \label{bre:bal:parseval-residual}
\end{equation}
La cola física se mide mediante \(B\widehat\Phi_k\), no sólo por
\(\widehat\Phi_k\): el campo depende de la combinación de canales
impuesta por \(B\), que puede tener núcleo. Se controlan las ecuaciones
retenidas y esta cola al aumentar \(H\). Si se emplea cuadratura de
Fourier en lugar de convoluciones exactas, se refina también la malla
angular para distinguir error de cuadratura y truncamiento.

Cuando \((\ii k\omega)^qI-A\) es invertible, el tamaño
\(\norm{[(\ii k\omega)^qI-A]^{-1}B\widehat\Phi_k}\) estima la
respuesta lineal al armónico omitido y detecta resonancias. No sustituye
un teorema de error para el problema no lineal, que requeriría controlar
la inversa del operador de balance y sus derivadas.

La secuencia de cálculo queda determinada: verificar la realización;
resolver fase, media y amplitud o el balance multiharmónico; reconstruir
la semilla con su fase; evaluar el residuo físico completo; continuar
hasta el campo original y analizar allí el PVI causal. Una raíz
algebraica aceptable debe superar tanto las condiciones de dominio
del resolvente como las comprobaciones de residuo y condicionamiento.
